\section{Conclusion et suite du projet}
\label{sec:conclusion}

\taub{} évalue l'agent sur ce qu'il \emph{fait}, l'état final de l'environnement, plutôt que sur ce qu'il \emph{dit}. Cette approche est objective et reproductible, et elle se prête bien à une évaluation automatique à grande échelle. Elle a pourtant une conséquence peu visible : en \texttt{airline}, le score mélange 24 tâches où il faut refuser et 26 où il faut agir, deux compétences que le chiffre agrégé ne distingue pas. Nos premiers essais montrent aussi qu'elle produit à la fois des faux négatifs (une demande satisfaite, mais un champ facultatif en trop) et des faux positifs (une violation de la politique sans trace dans la base). Le score dépend en outre de choix rarement explicités : modèle utilisateur, composition de \texttt{reward\_basis}, qui diffère selon le domaine, nombre d'essais. Enfin, plus d'un tiers des tâches \texttt{airline} et \texttt{retail} ont dû être corrigées. Un score \taub{} n'a donc de sens qu'accompagné de son protocole complet.

Notre premier essai sur l'effet du raisonnement l'illustre. Avec un budget de 4\,096 tokens, \texttt{Qwen3-14B} progresse sur toutes les mesures, et davantage sur la fiabilité ($\passk{2}$, $+18$\,\% relatif) que sur le taux de succès moyen. Mais ces écarts ne sont pas significatifs sur 50 tâches, et le score global ne dit pas si le raisonnement aide l'agent à mieux refuser ou à mieux agir. L'hypothèse reste donc ouverte. La tester correctement demande le protocole que nous construirons dans le rapport~2.

\paragraph{Pistes pour le rapport 2.} Nous retenons cinq questions de recherche candidates. Nous les préciserons et en choisirons deux ou trois selon le budget.
\begin{enumerate}[nosep]
  \item \textbf{Que révèle un score décomposé ?} On calcule le score séparément sur les tâches de refus et sur les tâches d'action, pour chaque agent évalué, et on regarde si des agents de même score global ont des profils différents (prudent ou complaisant). Cette analyse ne demande aucune simulation supplémentaire : elle s'applique aux trajectoires déjà obtenues, dont celles de \texttt{Qwen3-14B}. C'est l'analyse la moins coûteuse et la plus directement liée à la validité du score.
  \item \textbf{Le score mesure-t-il l'agent ou l'utilisateur simulé ?} À agent fixé, on fait varier le modèle utilisateur et on mesure la part de variance du score qui lui revient. La comparaison avec le mode \emph{solo} (\texttt{telecom}) isole de même le raisonnement de l'agent du bruit conversationnel.
  \item \textbf{Faux positifs et faux négatifs.} On annote à la main un échantillon de trajectoires pour estimer à quelle fréquence la note binaire contredit le respect réel de la politique. On compare ensuite avec une notation qui inclut les assertions en langue naturelle.
  \item \textbf{Combien d'essais pour conclure ?} On estime la variance entre essais et on en déduit le nombre de tâches et d'essais nécessaires pour séparer deux agents de manière significative. Les simulations \texttt{Qwen3-14B} déjà obtenues fournissent une première estimation.
  \item \textbf{Effet du raisonnement, de la taille et de la précision du modèle.} On prolonge l'essai de la section~\ref{sec:qwen} avec un protocole corrigé : utilisateur simulé fixe et distinct de l'agent, test apparié par tâche, score décomposé entre tâches de refus et tâches d'action. On fait varier le budget de raisonnement, la taille du modèle (14B contre 32B ou un modèle accessible par API) et sa quantification, pour savoir si le plafond observé tient au modèle ou à la formulation des tâches. On analyse aussi les traces de raisonnement, pour distinguer les vérifications préalables réussies des erreurs de calcul. Une extension au domaine \texttt{telecom} mesurerait enfin l'effet du raisonnement en \emph{dual-control}, où l'agent doit guider l'utilisateur.
\end{enumerate}
